% GENERATED FILE. Edit canonical lessons, then run npm run generate:books.
% canonical-source-hash: fnv1a64:4ea0062fa4b4ead6
% canonical-lessons: RU-C60-dog, RU-C60-cat, RU-C60-horse, RU-C60-cow, RU-C60-pig

\chapter{Dog, Cat, Horse, Cow, Pig}
\label{ch:ru-dog-cat-horse-cow-pig}
\hlchaptermodality{60}

\begin{chapteropening}
\textbf{By the end of this chapter:} \emph{I can name a dog, a cat, a horse, a cow and a pig.}

Complete the last lesson of chapter 60: i can name a dog, a cat, a horse, a cow and a pig.
\end{chapteropening}

\section[sobáka]{\ru{собака} (sobáka) — a dog}

\label{lesson:RU-C60-dog}



\begin{quote}
Before the new one: say the Russian for the weather, then the Russian for the world; peace.
\end{quote}

\subsection*{You'll want to know: \ru{собака}}
\textbf{\ru{собака}} (\emph{sobáka}) — \textquotedblleft{}a dog\textquotedblright{}.

\begin{grammarlens}[title={What you've built}]
One more word for this chapter.
\end{grammarlens}

\subsection*{Your turn}
\begin{itemize}
\raggedright
  \item \emph{Say it:} \emph{sobáka}
  \item \emph{Say it:} \emph{sobáka}, once more
  \item \emph{Say it:} say \emph{sobáka}
  \item \emph{Recall:} say the Russian for the weather, then the Russian for the world; peace, then say \emph{sobáka} again
\end{itemize}

\begin{tcolorbox}[breakable,colback=teal!4,colframe=teal!35!black,title={Before you move on}]
What does \ru{собака} mean? (\textquotedblleft{}A dog\textquotedblright{}.) Say it once more.
\end{tcolorbox}
\section[kóshka]{\ru{кошка} (kóshka) — a cat}

\label{lesson:RU-C60-cat}



\begin{quote}
Before the new one: say the Russian for the world; peace, then the Russian for a dog.
\end{quote}

\subsection*{You'll want to know: \ru{кошка}}
\textbf{\ru{кошка}} (\emph{kóshka}) — \textquotedblleft{}a cat\textquotedblright{}.

\begin{grammarlens}[title={What you've built}]
One more word for this chapter.
\end{grammarlens}

\subsection*{Your turn}
\begin{itemize}
\raggedright
  \item \emph{Say it:} \emph{kóshka}
  \item \emph{Say it:} \emph{kóshka}, once more
  \item \emph{Say it:} say \emph{kóshka}
  \item \emph{Recall:} say the Russian for the world; peace, then the Russian for a dog, then say \emph{kóshka} again
\end{itemize}

\begin{tcolorbox}[breakable,colback=teal!4,colframe=teal!35!black,title={Before you move on}]
What does \ru{кошка} mean? (\textquotedblleft{}A cat\textquotedblright{}.) Say it once more.
\end{tcolorbox}
\section[lóshad]{\ru{лошадь} (lóshad) — a horse}

\label{lesson:RU-C60-horse}



\begin{quote}
Before the new one: say the Russian for a dog, then the Russian for a cat.
\end{quote}

\subsection*{You'll want to know: \ru{лошадь}}
\textbf{\ru{лошадь}} (\emph{lóshad}) — \textquotedblleft{}a horse\textquotedblright{}.

\begin{grammarlens}[title={What you've built}]
One more word for this chapter.
\end{grammarlens}

\subsection*{Your turn}
\begin{itemize}
\raggedright
  \item \emph{Say it:} \emph{lóshad}
  \item \emph{Say it:} \emph{lóshad}, once more
  \item \emph{Say it:} say \emph{lóshad}
  \item \emph{Recall:} say the Russian for a dog, then the Russian for a cat, then say \emph{lóshad} again
\end{itemize}

\begin{tcolorbox}[breakable,colback=teal!4,colframe=teal!35!black,title={Before you move on}]
What does \ru{лошадь} mean? (\textquotedblleft{}A horse\textquotedblright{}.) Say it once more.
\end{tcolorbox}
\section[koróva]{\ru{корова} (koróva) — a cow}

\label{lesson:RU-C60-cow}



\begin{quote}
Before the new one: say the Russian for a cat, then the Russian for a horse.
\end{quote}

\subsection*{You'll want to know: \ru{корова}}
\textbf{\ru{корова}} (\emph{koróva}) — \textquotedblleft{}a cow\textquotedblright{}. It gives \textbf{\ru{молоко}}.

\begin{grammarlens}[title={What you've built}]
One more word for this chapter.
\end{grammarlens}

\subsection*{Your turn}
\begin{itemize}
\raggedright
  \item \emph{Say it:} \emph{koróva}
  \item \emph{Say it:} \emph{koróva}, once more
  \item \emph{Say it:} say \emph{koróva}
  \item \emph{Recall:} say the Russian for a cat, then the Russian for a horse, then say \emph{koróva} again
\end{itemize}

\begin{tcolorbox}[breakable,colback=teal!4,colframe=teal!35!black,title={Before you move on}]
What does \ru{корова} mean? (\textquotedblleft{}A cow\textquotedblright{}.) Say it once more.
\end{tcolorbox}
\section[svinyá]{\ru{свинья} (svinyá) — a pig}

\label{lesson:RU-C60-pig}



\begin{quote}
Before the new one: say the Russian for a horse, then the Russian for a cow.
\end{quote}

\subsection*{You'll want to know: \ru{свинья}}
\textbf{\ru{свинья}} (\emph{svinyá}) — \textquotedblleft{}a pig\textquotedblright{}.

\begin{grammarlens}[title={What you've built}]
One more word for this chapter.
\end{grammarlens}

\subsection*{Your turn}
\begin{itemize}
\raggedright
  \item \emph{Say it:} \emph{svinyá}
  \item \emph{Say it:} \emph{svinyá}, once more
  \item \emph{Say it:} say \emph{svinyá}
  \item \emph{Recall:} say the Russian for a horse, then the Russian for a cow, then say \emph{svinyá} again
\end{itemize}

\begin{tcolorbox}[breakable,colback=teal!4,colframe=teal!35!black,title={Before you move on}]
What does \ru{свинья} mean? (\textquotedblleft{}A pig\textquotedblright{}.) Say it once more.
\end{tcolorbox}
